CellCycleCells_dose_group_comparisons_density_by_dose_and_ploidy.png, CellCycleCells_dose_group_comparisons_sample_ecdf_by_dose_and_ploidy.png, CellCycleCells_dose_group_comparisons_sample_mean_pseudotime_by_dose.png: In the cell-cycle-associated tumor-cell population, pseudotime distributions differed significantly between untreated and gemcitabine-treated samples. A sample-aware ECDF permutation test showed a significant difference for 0 mg/kg versus combined 30 + 120 mg/kg treatment ($p = 0.0047$), and this remained significant after stratifying permutations by initial ploidy ($p = 0.0053$), indicating that the effect was not driven by imbalance between 2N and 4N samples. When doses were analyzed separately, both 30 mg/kg and 120 mg/kg differed from 0 mg/kg ($p = 0.034$ and $p = 0.018$, respectively), whereas 30 mg/kg and 120 mg/kg did not differ from each other ($p = 0.914$). Overall, gemcitabine treatment was associated with a reproducible shift in pseudotime distribution relative to untreated controls, without clear evidence of a dose-dependent separation between the two treated groups. 

CellCycleCells_pseudotime_shift_associations_TGI_AUC_ecdf_distance_vs_TGI_AUC.png: Among gemcitabine-treated samples, larger pseudotime shifts from the ploidy-matched untreated reference were associated with higher TGI AUC by Pearson correlation (ECDF RMSE; $r = 0.71$, $p = 0.049$), although the corresponding Spearman correlation was not significant ($\rho = 0.64$, $p = 0.086$), suggesting a positive but sample-size-sensitive trend. Pseudotime-shift magnitude did not associate with mean cell ploidy ($p = 0.65$; treated-only $p = 0.46$) or initial ploidy ($p = 0.64$; treated-only $p = 0.47$).
